\section{Introduction}
\label{sec:introduction}

\subsection{Empathy as a field: a hundred years of scattered contacts}

From the early psychoanalysts \citep{ferenczi1933, klein1946} through contemporary cognitive neuroscience \citep{decety2011, singer2014}, the psychology of empathy has accumulated several major, internally rich, but poorly articulated lines of inquiry. The psychoanalytic tradition describes splitting, repression, and the defensive mechanisms that turn the experience of the other into its opposite \citep{freud1915, fairbairn1952, bion1962, fonagy2002}. Attachment theory develops environmental antecedents---the quality of early relationships as a predictor of adult affective resonance \citep{bowlby1969, main1986, lyonsruth1999}. Cognitive neuroscience separates cognitive empathy (theory of mind) from affective empathy (resonance), pointing to distinct neural substrates \citep{shamaytsoory2011, decety2011}. The mirror-neuron literature proposes a biological mechanism of affective resonance \citep{rizzolatti2004, iacoboni2009}. The clinical tradition of psychopathy, from \citet{cleckley1941} through \citet{hare1991} and \citet{blair2013}, treats affective deficit as a structural property of a specific subpopulation. Contemplative traditions, from the Buddha to contemporary nondual literature \citep{godman1985, dunne2011, josipovic2014}, describe a distinct stratum of consciousness---a witnessing stance, meta-awareness---that, according to these traditions, modulates all other psychological functions.

Each of these lines has developed independently, with its own instruments, its own vocabulary, its own validation standards. This is not a theoretical reproach: detailed work within each paradigm yields knowledge that would be impossible under eclectic merger. But it is an empirical problem: the researcher who wishes to measure empathy in a concrete person and to make warranted claims about that person's resources and vulnerabilities finds themselves with five or six partial maps of one territory, none of which covers the whole.

This paper does not propose a seventh map on top of the six existing ones. It proposes a \emph{measurement framework} in which the constructs already measurable within each tradition receive a shared quantitative coordinate system that permits comparison and cross-reference within a single sample. We call this framework \DEIC{} (\textit{Dynamic Empathy and Inner Conflict}---a working name emphasizing the interaction between empathy and defensive mechanisms). Psychometrically, \DEIC{} is a six-factor model with one additional formative construct (attention, A), evaluated by confirmatory analysis on a sample of $N = 1426$.

\subsection{Empathy as default: an inverted measurement frame}
\label{subsec:empathy_as_default}

Before describing the specific factors of the model, we state the inverting thesis on which \DEIC{} is built and which distinguishes it from most existing empathy instruments. The standard measurement tradition---from IRI \citep[Interpersonal Reactivity Index;][]{davis1983} through the Empathy Quotient \citep{baroncohen2011} to contemporary multidimensional instruments---proceeds from the implicit premise that empathy is an individual ability varying in magnitude: some people have more of it, others less, and the task of the instrument is to quantify that magnitude. We proceed from the opposite premise.

\textbf{Empathy is the default state of the human psyche.} Affective resonance with another's state is not an achievement, not a resource, not a trait that some have developed and others have not; it is a basic function of the primate social brain \citep{rizzolatti2004, iacoboni2009, decety2011}, neurologically supported by mirror-neuron and affective-resonance systems that operate in infants from the first months of life. \emph{Observed between-person differences in empathy are not differences in its quantity, but differences in the mechanisms that suppress it.} The model therefore proposes to measure not empathy as such, but \emph{the forms and strength of non-empathy}---the active and passive processes that block the underlying capacity.

From this inversion follows a distinction that is the central conceptual instrument of \DEIC{}. There are two architecturally different channels of non-empathy.

\textbf{Passive blocking.} The empathic response begins but is blocked by the defensive system, because the experience itself---\emph{regardless of whether it is another's or one's own}---is intolerable to the psyche in its current configuration. One common oversimplification must be set aside immediately: passive blocking is not a block on ``feeling the other'' separate from ``feeling oneself.'' It is one and the same channel. The capacity to bear affective resonance with another's pain and the capacity to bear one's own inner material are structurally a single function, not two separate ones: the reversible sensitivity of a containing system that either holds affect or does not. The psychoanalytic and mentalization traditions record this explicitly: \citet{bion1962}'s $\alpha$-function---the capacity to metabolize emotional experience---operates symmetrically on one's own and received states; \citet{fonagy2002}'s reflective function is by definition bidirectional (the capacity to mentalize others is predicted by the capacity to mentalize oneself, and vice versa); \citet{winnicott1965}'s holding is possible in the mother precisely to the extent that she can bear her own states. The neurobiological line of \citet{schore2001} shows that the substrate of self-regulation and other-attunement is shared: the right-hemisphere networks responsible for regulating one's own arousal are the same networks that enable affective attunement to another. A person in the regime of passive blocking therefore \emph{cannot} feel in both directions simultaneously: they are equally cut off from tears over another's loss and from their own fatigue, equally blind to rising anxiety in a partner and in themselves. Phenomenologically this appears as numbing, dissociation, affective flattening, an inner conflict of ``I ought to feel but I cannot''---and the ``ought'' applies both to response to the other and to contact with one's own content. The psychometric coordinate is the factor ID (Internal Dissonance). Historical descriptions: \citet{ferenczi1933} on ``identification with the aggressor,'' \citet{klein1946} on splitting, \citet{bowlby1969} and \citet{main1986} on disorganized attachment models, and the clinical dissociative literature \citep{putnam1997, vanderkolk2014}.

\textbf{Active switching-off.} The empathic channel is intact and functionally available, but it does not fire automatically; it engages or disengages depending on contextual suitability and a moral narrative about the situation. A person in this regime \emph{can} feel selectively and can \emph{refrain} from feeling when it is convenient. Phenomenologically this is not numbness but control---smooth switchability. The psychometric coordinate is the factor P (psychopathic pattern). Historical descriptions: \citet{cleckley1941}, \citet{karpman1941}, the clinical tradition \citep{hare1991, lilienfeld1990}, and the neuroscientific line of Blair \citep{blair1995, blair2013}.

These two channels are independent. They are \emph{not} ``the same thing in different proportions'' and do not form a single spectrum. A person may have passive blocking in the complete absence of any psychopathic narrative (the typical clinical profile of depression and complex trauma); may have the conditional psychopathic architecture in the absence of ID (the primary pattern of \citet{karpman1941}, empirically reproduced in our sample at $n = 144$); or may have both at once (the secondary pattern, $n = 65$, $\mathrm{ID} \approx +0.90$ and $\mathrm{CA} \approx +0.98$). Their independence is one of the principal empirical results of this paper (§\ref{subsec:architecture}), and from it follow substantially different clinical strategies.

This inversion of frame is supported by two external empirical anchors that we take to be the strongest arguments in its favor. First, the meta-analytic estimate of \citet{depow2025}: trait measures of empathy explain only 3--15\,\% of the variance in empathic experience in everyday life. If empathy were a stable individual trait that some have and others do not, standard trait instruments would capture a substantially larger share of the variance. The observed 3--15\,\% suggests that the empathic response is a state strongly dependent on situational factors, among which the \emph{activation or deactivation of blocking mechanisms} plays the central role, rather than fluctuations in a baseline resource. Second, the result of \citet{ogawa1997}: absence of parental responsiveness in infancy predicts dissociative symptoms nineteen years later with an effect size of roughly 50\,\% of explained variance---stronger than the prediction from traumatic events themselves. This is consistent with \citet{fonagy2002}'s hypothesis: not the stressful event as such, but the absence of a mentalizing field within which the event could be metabolized, forms the structural dissociative defense. That is, the blocking mechanism forms precisely where the adult capacity to co-experience a child's inner states fails or never arises.

A third anchor is structural. \citet{digiuseppe2024}, in a network analysis of defensive mechanisms, showed that sixteen different defenses form a densely connected integrated system (not independent traits), in which self-assertion serves as the central node and the remaining defenses as peripheral nodes carrying different functional loads. In our data, correlations among subscales of passive blocking (dissociation, suppression, affective isolation, flattening, internal dissonance, masking) lie in the range $r = 0.77$--$0.94$. This is not an instrument-level redundancy artifact but an independent reproduction of what Di Giuseppe's network analysis found on its own: defensive mechanisms are \emph{one integrated functional system}, not a set of parallel individual traits. This is why in our 6-factor CFA these subscales collapse into a single latent factor ID---they are not psychometrically discriminable not because the instrument fails to distinguish them, but because in reality they are different phenomenological expressions of one suppression system.

These three anchors---the low explanatory power of trait empathy, the strong prediction of dissociation through the absence of a mentalizing response, and the network integration of defenses---converge independently on a single conclusion: the central object of measurement is \emph{the mechanisms of non-empathy}, not empathy as a trait. The \DEIC{} model is built on this inversion, and all its subsequent claims should be read through that frame.

\subsection{Five points where traditions have touched without meeting}

The reason the \DEIC{} architecture has a chance of being useful lies not in its ambition but in its narrowness. We have selected for modeling those constructs in which different traditions empirically describe \emph{very similar} phenomena but have lacked a shared instrument to confirm it.

\paragraph{Internal dissonance (ID).}
In the psychoanalytic tradition---splitting and ego splitting-off \citep{ferenczi1933, klein1946}; in attachment theory---disorganized internal working models \citep{main1986}; in the clinical dissociative literature---fragmentation in response to trauma \citep{putnam1997, vanderkolk2014}. All three describe the same phenomenon: a person's inner states bear the mark of active blocking, regulation, defensive work, rather than merely content. Psychometrically this should be a single factor if the instrument is correctly constructed. We show that it is.

\paragraph{Childhood alienation (CA).}
In attachment theory---the absence of a secure base \citep{bowlby1969, ainsworth1978}; in mentalization theory---the absence of a mentalizing response \citep{fonagy2002}; in the clinical ACE literature---stressful childhood events \citep{felitti1998}. CA is the environmental condition in which a child's inner states do not meet an adequate response. We instrument it narrowly (not the full ACE index) because, per Fonagy's hypotheses, it is the \emph{absence of a mentalizing field}, not the occurrence of stressful events as such, that predicts downstream dysregulation. The substantive meaning of this narrow formulation: a child learns to distinguish and hold their own states through an adult capable of seeing and naming those states. If no such adult is available, the child does not lose ``information about others'' separately from ``information about the self''; the child fails to develop the general capacity to metabolize affective experience as such. This is why CA $\to$ ID in our model is not ``trauma reduced empathy toward others''; it is ``the absence of a mentalizing mirror did not permit the capacity to bear affective material in both directions at once to form.'' The clinical signature of an adult with high CA and high ID is uniform impermeability both outward and inward: such a person has access neither to their own sadness nor to a partner's sadness, because the only available mechanism is an automatic joint block, not a selective closure of one direction.

\paragraph{Psychopathy as a moral signature (P).}
The clinical tradition \citep{cleckley1941, hare1991} describes psychopathy in the frame of affective deficit and the instrumental use of others; Blair's neuroscience \citep{blair1995, blair2013}---through reduced amygdalar response to another's distress; evolutionary psychology \citep{mealey1995, lalumiere2001}---as a low-frequency stable strategy. We deliberately do not reproduce the deficit frame. In \DEIC{}, the psychopathic pattern is decomposed into two layers: the \textbf{architectural} (conditional rather than unconditional affective coupling---a morally neutral variation, functional in several social roles; see §\ref{sec:discussion}.5) and the \textbf{narrative} (P---a stable assertion ``this is allowed for me,'' determining how the architecture is used socially). The factor P models the second layer. This two-layer shift---from ``broken function'' to ``architecture plus narrative''---alters the clinical, legal, and political implications and is one of the key contributions of the work.

\paragraph{Cognitive vs. affective empathy (\Ecog{} vs. \EAF{}).}
In neuropsychology---two systems with distinguishable substrates \citep{shamaytsoory2011, decety2011}; in mentalization theory---mentalization vs. embodied empathy \citep{fonagy2002}; in clinical work---the key distinction for understanding why an ``intelligent psychopath'' can be accurate in predicting behavior yet indifferent to suffering. We show empirically that these two constructs are separable and differently related to P. An everyday illustration of this distinction is one of the most structurally informative observations usually ignored in trait-measurement literature: the same person may weep over the fate of a literary or film character and remain cold to the suffering of a concrete real person next to them. This is not a contradiction and not ``hypocrisy''---it is a direct signature of channel splitting. Identification with a narrative (a fictional character within a safe storytelling frame) engages primarily the \Ecog{} line: mentalization of a point of view without automatic affective coupling to a real, threatening, unsafe presence of an other. A real suffering person calls for \EAF{} in its full form---affective resonance in a situation that is not governed by a narrative frame and that can make demands. When \EAF{} is blocked by the defensive structure (ID) or regulated conditionally (P), the \Ecog{} channel remains partially or fully intact, and the person experiences a state in which fictional suffering moves them but real suffering does not. In the DEIC architecture this is not an anomaly but a predictable consequence of the separability of the two channels.

\paragraph{Attention as witnessing (A).}
In the contemplative tradition---\textit{satipa\d{t}\d{t}h\=ana} (Buddhism), \textit{s\=ak\d{s}\=i} (Advaita Ved\=anta), \textit{rigpa} (Dzogchen), ``self-remembering'' \citep[Gurdjieff through][]{ouspensky1949}; in the psychoanalytic tradition---the observing ego \citep{sterba1934, kris1956}; in contemporary mindfulness science---meta-awareness \citep{dunne2011, josipovic2014}; in mentalization theory---the reflective function \citep{fonagy1997}. Each of these lines describes a structurally distinct stratum of consciousness: not the content of the psyche but a \emph{relation of observation} to that content. We return to A repeatedly---it is the most theoretically rich and simultaneously the most difficult to instrument of the factors in the model.

\medskip
An important qualification: we do not claim that \DEIC{} describes these traditions \emph{in full} or that it is the \emph{first} to catch the structure they describe. Each of the literatures listed above has measured its construct at greater depth than we do. Our contribution lies in having measured \emph{all of the listed constructs in one sample, with one instrument, under shared methodological standards}. This makes it possible, for the first time, to test their mutual relations not at the level of theoretical reconstruction but at the level of path-model coefficients.

\subsection{Balkanization as failure mode, not as discipline}

Before turning to method, we state the position openly. The tradition of holding that ``a narrow frame = strict science,'' applied in twentieth-century psychology, has resulted in a situation where incompatible microtheories coexist as the norm and are called a discipline. The theory of narcissism is not coordinated with the theory of psychopathy; neither is coordinated with attachment theory; cognitive empathy is measured with instruments architecturally incompatible with affective empathy. This is precisely what \citet{meehl1978} diagnosed half a century ago as the failure mode of psychological theorizing. Physics, chemistry, and biology demand coherence across scales; psychology has allowed itself to abandon that demand under the rhetoric of ``caution.''

We do not adopt that caution as a virtue. We hold that: if the data contain constructs described by different traditions as the same phenomenon under different names, psychometrics should demonstrate so. If the \DEIC{} architecture makes predictions, it is obliged to state them in falsifiable form---in that manner, not through a ``narrow restriction of scope'' that exempts from testing what the model actually predicts. This is the methodological wager of the paper: ambition in content, discipline in how the question is posed.

\subsection{Positioning: three strata, one sentence}

Relative to three reference planes the work stands as follows. In relation to the \textbf{contemplative tradition} (Buddhism, Advaita, Dzogchen, Hesychasm, Gurdjieff, psychoanalysis at its contemplative edge), \DEIC{} stands in the position of \emph{student-translator}: the constructs these traditions distinguished without psychometric instruments for millennia receive coordinates in our framework---and our contribution lies precisely there, not in a ``new integration.'' In relation to \textbf{academic psychology}, \DEIC{} stands in a \emph{confrontational} position: balkanization is not a discipline but a pathology of the field, and a framework of shared referents for what different schools describe in incompatible languages is something the field should have but does not. In relation to the \textbf{integral tradition} (Wilber and AQAL, structural precedents), \DEIC{} stands in a \emph{structurally inverse} position: not a master map onto which traditions land as special cases, but an operational dialect of those same traditions, deployed against fragmentation.

One sentence: \DEIC{} is an operational dialect of what the contemplative tradition measured without instruments, deployed against the balkanization of academic psychology. Student to the tradition, challenger to the field.

\subsection{What \DEIC{} does not claim}

Because the program is positioned as a framework crossing several traditions, we delimit its scope at the outset. \DEIC{} does \textbf{not} claim:
\begin{itemize}
    \item To be the first or the only unifying framework for empathy. Unification is not our verb.
    \item To contain a final operationalization of A, witnessing, or mindfulness. A in its current form is a formative composite, not a reflective latent; this limitation is empirically established (§\ref{sec:results}.7) and is to be corrected in the following wave.
    \item To replace any of the traditions we cite. Psychoanalysis, attachment theory, the cognitive neuroscience of empathy, the mirror-neuron literature, the clinical tradition of psychopathy, and the contemplative schools each possess depth and specificity within their own paradigms that we do not reproduce.
    \item To offer a theory of consciousness or an explanation of the hard problem \citep{chalmers1995}. We measure \emph{functions} within consciousness; we assert nothing about the nature of subjective experience as such.
    \item To deliver clinical recommendations on the basis of a single sample. Everything clinically prescriptive in the Discussion is formulated as \emph{hypotheses for testing}, not as immediate protocols.
\end{itemize}

These limits are not a reluctant retreat but a structural condition for producing work that does not collapse into a meta-theoretical claim. The program is built out of many narrow, falsifiable studies; this paper is one of them.

\subsection{Structure of the paper}

Section~\ref{sec:methods} describes the DEIC-Inventory, the operationalization of the six factors through a 6-factor oblique CFA, the construction of A as a formative composite, and the statistical approach (conditional process SEM, parallel mediation, unified SEM, A-CFA). Section~\ref{sec:results} presents the factor structure, partial correlations, the central conditional-process analysis, parallel mediation, a unified SEM carrying all paths simultaneously, the empirical confirmation of the integrated survivor, the topological distinction between primary and secondary psychopathy, and the empirical demonstration of the formative nature of A. Section~\ref{sec:discussion} takes up the implications of the architecture, isolates A as an ontologically asymmetric factor (field-condition, not field-content), formulates falsifiable predictions and points of divergence with existing theories, and situates the paper within a future program. Limitations and directions for future work appear in Sections~\ref{sec:limitations} and~\ref{sec:future}.
